\documentclass[12pt]{letter}
\usepackage[margin=1in]{geometry}
\usepackage{microtype}
\usepackage{hyperref}

\signature{Jacob K. Moutouama and\\
Tom E.X. Miller\\
On behalf of all co-authors}

\address{}
\date{\today}

\begin{document}

\begin{letter}{%
  Dr.\ Peter H. Thrall\\
  Editor-in-Chief, \textit{Ecology Letters}\\
  CSIRO National Collections \& Marine Infrastructure\\
  GPO Box 1600, Canberra ACT 2601, Australia
}

\opening{Dear Dr.\ Thrall,}

We are pleased to submit our manuscript, \textit{``When mutualism weakens: aridity and herbivory shape endophyte benefits in cool-season grasses''}, for consideration as a Letter in \textit{Ecology Letters}.

The stress-gradient hypothesis predicts that mutualistic benefits increase under environmental stress, yet the conditions under which stress enhances versus undermines symbiosis remain poorly understood. We tested this using field experiments along a natural precipitation gradient with herbivore exclusion treatments, combined with a greenhouse experiment manipulating soil origin, to examine how aridity and herbivory shape the demographic effects of vertically transmitted fungal endophytes (\textit{Epichlö} spp.) in three native cool-season grasses.

Contrary to predictions, endophyte benefits did not increase under aridity; survival and growth benefits were strongest under  moist conditions and weakened toward dry range edges. Aridity primarily drove effects on survival and growth, while herbivory shaped reproductive traits. Soil origin did not modify outcomes, ruling out precipitation-driven soil properties as a driver. These findings show that endophytes can become liabilities under stress, with implications for plant-fungal mutualisms under global change.

This work offers a direct empirical test of a foundational prediction using a replicated natural gradient, tracks multiple vital rates across seven sites and two herbivory treatments, and reveals a counterintuitive pattern with broad relevance to predicting the fate of plant--microbe interactions as climate change intensifies drought and alters herbivore pressure.

We have opted out of double-blind review because this paper builds upon our previous work, so our identities could be easily deduced. Our identities are also evident from our GitHub repository, where we direct reviewers to find all data and code necessary to reproduce our analyses.


\closing{Sincerely yours,}

\end{letter}
\end{document}
